% ==============================================================================
% LAMPIRAN-LAMPIRAN SKRIPSI
% Sesuai Pedoman Penulisan Skripsi FST UNIBA Madura 2026
% ==============================================================================

\clearpage
\phantomsection
\chapter*{LAMPIRAN}
\addcontentsline{toc}{chapter}{LAMPIRAN}

\vspace{0.5cm}

% ==============================================================================
% LAMPIRAN 1: ORKESTRASI DOCKER COMPOSE
% ==============================================================================
\section*{Lampiran 1: Berkas Orkestrasi Kluster Kontainer Heterogen (Docker Compose)}

Berikut adalah berkas orkestrasi \texttt{docker-compose.yml} yang mendefinisikan pembatasan kuota sumber daya cgroups v2 (\textit{vCPU} dan batas memori RAM) dengan rasio kapasitas heterogen 1:2:4, isolasi jaringan privat \textit{Docker Bridge Network}, profil peladen penyeimbang beban NGINX dan HAProxy, serta instrumen telemetri cAdvisor:

\begin{lstlisting}[language=bash, caption={Berkas docker-compose.yml Kluster Kontainer Heterogen}]
version: '3.8'

# Topologi Heterogen: Rasio cgroups 1:2:4
networks:
  skripsi-net:
    driver: bridge
    ipam:
      config:
        - subnet: 172.28.0.0/16

services:
  # Backend 1: Tier Small (0.5 vCPU, 256 MB RAM) -> Rasio 1
  web-backend-1:
    build:
      context: ./backend
      dockerfile: Dockerfile
    container_name: web-backend-1
    hostname: web-backend-1
    restart: unless-stopped
    environment:
      - PORT=8080
    deploy:
      resources:
        limits:
          cpus: '0.50'
          memory: 256M
    cpus: 0.50
    mem_limit: 256m
    networks:
      - skripsi-net

  # Backend 2: Tier Medium (1.0 vCPU, 512 MB RAM) -> Rasio 2
  web-backend-2:
    build:
      context: ./backend
      dockerfile: Dockerfile
    container_name: web-backend-2
    hostname: web-backend-2
    restart: unless-stopped
    environment:
      - PORT=8080
    deploy:
      resources:
        limits:
          cpus: '1.00'
          memory: 512M
    cpus: 1.00
    mem_limit: 512m
    networks:
      - skripsi-net

  # Backend 3: Tier Large (2.0 vCPU, 1024 MB RAM) -> Rasio 4
  web-backend-3:
    build:
      context: ./backend
      dockerfile: Dockerfile
    container_name: web-backend-3
    hostname: web-backend-3
    restart: unless-stopped
    environment:
      - PORT=8080
    deploy:
      resources:
        limits:
          cpus: '2.00'
          memory: 1024M
    cpus: 2.00
    mem_limit: 1024m
    networks:
      - skripsi-net

  # Load Balancer NGINX (Default Active Profile)
  loadbalancer-nginx:
    image: nginx:1.27-alpine
    container_name: lb-nginx
    restart: unless-stopped
    ports:
      - "80:80"
    volumes:
      - ${NGINX_CONF:-./loadbalancer/nginx/nginx-wrr.conf}:/etc/nginx/nginx.conf:ro
      - ./logs/nginx:/var/log/nginx
    depends_on:
      - web-backend-1
      - web-backend-2
      - web-backend-3
    networks:
      - skripsi-net
    profiles:
      - nginx
      - default

  # Load Balancer HAProxy (Alternative Profile)
  loadbalancer-haproxy:
    image: haproxy:2.8-alpine
    container_name: lb-haproxy
    restart: unless-stopped
    ports:
      - "80:80"
      - "8404:8404" # Stats Dashboard
    volumes:
      - ${HAPROXY_CONF:-./loadbalancer/haproxy/haproxy-wrr.cfg}:/usr/local/etc/haproxy/haproxy.cfg:ro
      - ./logs/haproxy:/var/log
    depends_on:
      - web-backend-1
      - web-backend-2
      - web-backend-3
    networks:
      - skripsi-net
    profiles:
      - haproxy

  # Telemetri cAdvisor (Monitoring cgroups real-time)
  cadvisor:
    image: gcr.io/cadvisor/cadvisor:v0.49.1
    container_name: skripsi-cadvisor
    restart: unless-stopped
    privileged: true
    volumes:
      - /:/rootfs:ro
      - /var/run:/var/run:ro
      - /sys:/sys:ro
      - /var/lib/docker/:/var/lib/docker:ro
      - /dev/disk/:/dev/disk:ro
    ports:
      - "8081:8080"
    networks:
      - skripsi-net
    profiles:
      - telemetry
\end{lstlisting}

\newpage

% ==============================================================================
% LAMPIRAN 2: KONFIGURASI NGINX LOAD BALANCER
% ==============================================================================
\section*{Lampiran 2: Berkas Konfigurasi Mesin Penyeimbang Beban NGINX}

Berikut adalah berkas konfigurasi NGINX untuk algoritma \textit{Weighted Round Robin} (\texttt{nginx-wrr.conf}) dan \textit{Weighted Least Connections} (\texttt{nginx-wlc.conf}) yang mengintegrasikan pencatatan log akses berformat kustom untuk keperluan audit deviasi distribusi beban:

\begin{lstlisting}[language=bash, caption={Konfigurasi nginx-wrr.conf (Weighted Round Robin)}]
user nginx;
worker_processes auto;
error_log /var/log/nginx/error.log warn;
pid /var/run/nginx.pid;

events {
    worker_connections 10240;
    use epoll;
    multi_accept on;
}

http {
    include /etc/nginx/mime.types;
    default_type application/octet-stream;

    log_format audit_log '$remote_addr - [$time_local] "$request" '
                         '$status $body_bytes_sent rt=$request_time '
                         'upstream_addr="$upstream_addr" '
                         'upstream_status="$upstream_status" '
                         'upstream_response_time="$upstream_response_time"';

    access_log /var/log/nginx/access.log audit_log buffer=64k flush=1s;

    sendfile on;
    tcp_nopush on;
    tcp_nodelay on;
    keepalive_timeout 65;

    # Upstream: Weighted Round Robin (WRR - Bobot 1:2:4 sesuai kapasitas cgroups)
    upstream backend_cluster {
        server web-backend-1:8080 weight=1 max_fails=3 fail_timeout=10s;
        server web-backend-2:8080 weight=2 max_fails=3 fail_timeout=10s;
        server web-backend-3:8080 weight=4 max_fails=3 fail_timeout=10s;
        keepalive 64;
    }

    server {
        listen 80;
        server_name localhost;

        location / {
            proxy_pass http://backend_cluster;
            proxy_http_version 1.1;
            proxy_set_header Connection "";
            proxy_set_header Host $host;
            proxy_set_header X-Real-IP $remote_addr;
            proxy_set_header X-Forwarded-For $proxy_add_x_forwarded_for;
            proxy_next_upstream error timeout invalid_header http_500 http_502 http_503 http_504;
        }

        location /nginx_status {
            stub_status on;
            access_log off;
        }
    }
}
\end{lstlisting}

\vspace{0.4cm}

\begin{lstlisting}[language=bash, caption={Konfigurasi nginx-wlc.conf (Weighted Least Connections)}]
# Upstream blok untuk Weighted Least Connections (least_conn)
upstream backend_cluster {
    least_conn;
    server web-backend-1:8080 weight=1 max_fails=3 fail_timeout=10s;
    server web-backend-2:8080 weight=2 max_fails=3 fail_timeout=10s;
    server web-backend-3:8080 weight=4 max_fails=3 fail_timeout=10s;
    keepalive 64;
}
\end{lstlisting}

\newpage

% ==============================================================================
% LAMPIRAN 3: KONFIGURASI HAPROXY LOAD BALANCER
% ==============================================================================
\section*{Lampiran 3: Berkas Konfigurasi Mesin Penyeimbang Beban HAProxy}

Berikut adalah berkas konfigurasi HAProxy untuk algoritma \textit{Weighted Round Robin} (\texttt{haproxy-wrr.cfg}) dan \textit{Weighted Least Connections} (\texttt{haproxy-wlc.cfg}):

\begin{lstlisting}[language=bash, caption={Konfigurasi haproxy-wrr.cfg (Weighted Round Robin)}]
global
    log stdout format raw local0
    maxconn 20000
    stats socket /tmp/haproxy.sock mode 660 level admin

defaults
    log     global
    mode    http
    option  httplog
    option  dontlognull
    retries 3
    timeout connect 5000ms
    timeout client  50000ms
    timeout server  50000ms

frontend http_in
    bind *:80
    default_backend web_backends

backend web_backends
    balance roundrobin
    option httpchk GET /health
    http-check expect status 200
    # Alokasi bobot 1:2:4 sesuai perbandingan cgroups
    server web1 web-backend-1:8080 check weight 1
    server web2 web-backend-2:8080 check weight 2
    server web3 web-backend-3:8080 check weight 4

listen stats
    bind *:8404
    stats enable
    stats uri /stats
    stats refresh 5s
\end{lstlisting}

\vspace{0.4cm}

\begin{lstlisting}[language=bash, caption={Konfigurasi haproxy-wlc.cfg (Weighted Least Connections)}]
backend web_backends
    balance leastconn
    option httpchk GET /health
    http-check expect status 200
    # Alokasi bobot 1:2:4 dengan algoritma leastconn
    server web1 web-backend-1:8080 check weight 1
    server web2 web-backend-2:8080 check weight 2
    server web3 web-backend-3:8080 check weight 4
\end{lstlisting}

\newpage

% ==============================================================================
% LAMPIRAN 4: KODE SUMBER BACKEND SERVICE GOLANG & DOCKERFILE
% ==============================================================================
\section*{Lampiran 4: Implementasi Layanan Backend Golang \& Berkas Citra Docker}

Berikut adalah implementasi peladen \textit{backend} berbasis bahasa pemrograman Go (\texttt{main.go}) yang mengekspos \textit{endpoint} intensif \texttt{/api/cpu} dan antrean soket \texttt{/api/io}, serta berkas \texttt{Dockerfile} berbasis kompilasi multi-tahap (\textit{multi-stage build}):

\begin{lstlisting}[language=C, caption={Kode Sumber Backend Service (main.go)}]
package main

import (
	"crypto/sha256"
	"encoding/hex"
	"encoding/json"
	"fmt"
	"log"
	"net/http"
	"os"
	"strconv"
	"time"
)

type Response struct {
	Status        string `json:"status"`
	ContainerID   string `json:"container_id"`
	Endpoint      string `json:"endpoint"`
	ExecutionTime string `json:"execution_time_ms"`
	Timestamp     int64  `json:"timestamp"`
	Result        string `json:"result,omitempty"`
}

var containerID string

func init() {
	var err error
	containerID, err = os.Hostname()
	if err != nil {
		containerID = "unknown-container"
	}
}

func healthHandler(w http.ResponseWriter, r *http.Request) {
	w.Header().Set("Content-Type", "application/json")
	w.Header().Set("X-Backend-Server", containerID)
	json.NewEncoder(w).Encode(Response{
		Status:      "healthy",
		ContainerID: containerID,
		Endpoint:    "/health",
		Timestamp:   time.Now().UnixMilli(),
	})
}

// cpuHandler menjalankan hash SHA-256 berulang untuk memicu komputasi vCPU
func cpuHandler(w http.ResponseWriter, r *http.Request) {
	start := time.Now()
	iterations := 10000
	if iterParam := r.URL.Query().Get("iter"); iterParam != "" {
		if val, err := strconv.Atoi(iterParam); err == nil && val > 0 {
			iterations = val
		}
	}

	data := []byte("uniba-madura-skripsi-benchmarking-seed-" + strconv.FormatInt(start.UnixNano(), 10))
	for i := 0; i < iterations; i++ {
		hash := sha256.Sum256(data)
		data = hash[:]
	}

	elapsed := float64(time.Since(start).Microseconds()) / 1000.0

	w.Header().Set("Content-Type", "application/json")
	w.Header().Set("X-Backend-Server", containerID)
	json.NewEncoder(w).Encode(Response{
		Status:        "success",
		ContainerID:   containerID,
		Endpoint:      "/api/cpu",
		ExecutionTime: fmt.Sprintf("%.2f", elapsed),
		Timestamp:     time.Now().UnixMilli(),
		Result:        hex.EncodeToString(data),
	})
}

// ioHandler mensimulasikan latensi operasi basis data (150ms non-blocking)
func ioHandler(w http.ResponseWriter, r *http.Request) {
	start := time.Now()
	delay := 150 * time.Millisecond
	if delayParam := r.URL.Query().Get("delay"); delayParam != "" {
		if ms, err := strconv.Atoi(delayParam); err == nil && ms >= 0 {
			delay = time.Duration(ms) * time.Millisecond
		}
	}

	time.Sleep(delay)
	elapsed := float64(time.Since(start).Microseconds()) / 1000.0

	w.Header().Set("Content-Type", "application/json")
	w.Header().Set("X-Backend-Server", containerID)
	json.NewEncoder(w).Encode(Response{
		Status:        "success",
		ContainerID:   containerID,
		Endpoint:      "/api/io",
		ExecutionTime: fmt.Sprintf("%.2f", elapsed),
		Timestamp:     time.Now().UnixMilli(),
		Result:        fmt.Sprintf("I/O delay of %v completed", delay),
	})
}

func main() {
	port := os.Getenv("PORT")
	if port == "" {
		port = "8080"
	}

	mux := http.NewServeMux()
	mux.HandleFunc("/", healthHandler)
	mux.HandleFunc("/health", healthHandler)
	mux.HandleFunc("/api/cpu", cpuHandler)
	mux.HandleFunc("/api/io", ioHandler)

	server := &http.Server{
		Addr:         ":" + port,
		Handler:      mux,
		ReadTimeout:  15 * time.Second,
		WriteTimeout: 15 * time.Second,
		IdleTimeout:  60 * time.Second,
	}

	log.Printf("[INFO] Backend container %s listening on port %s...", containerID, port)
	if err := server.ListenAndServe(); err != nil && err != http.ErrServerClosed {
		log.Fatalf("[FATAL] Server error: %v", err)
	}
}
\end{lstlisting}

\vspace{0.4cm}

\begin{lstlisting}[language=bash, caption={Berkas Multi-Stage Dockerfile Backend}]
# Tahap 1: Kompilasi Binary Go
FROM golang:1.22-alpine AS builder
WORKDIR /app
COPY go.mod ./
COPY main.go ./
RUN CGO_ENABLED=0 GOOS=linux go build -ldflags="-s -w" -o backend-app main.go

# Tahap 2: Runtime Image Ringan
FROM alpine:3.20
RUN apk --no-cache add ca-certificates tzdata
WORKDIR /root/
COPY --from=builder /app/backend-app .
ENV PORT=8080
EXPOSE 8080
CMD ["./backend-app"]
\end{lstlisting}

\newpage

% ==============================================================================
% LAMPIRAN 5: SKRIP UJI BEBAN GRAFANA K6
% ==============================================================================
\section*{Lampiran 5: Skrip Skenario Pengujian Beban Grafana k6}

Berikut adalah skrip otomasi pengujian beban stres k6 untuk skenario \textit{CPU-bound} (\texttt{cpu-test.js}) dan \textit{I/O-bound} (\texttt{io-test.js}) dengan penambahan pengguna virtual bertahap (100 $\rightarrow$ 500 $\rightarrow$ 1.000 VUs):

\begin{lstlisting}[language=bash, caption={Skrip Skenario CPU-bound (cpu-test.js)}]
import http from 'k6/http';
import { check, sleep } from 'k6';
import { Counter, Rate, Trend } from 'k6/metrics';

export const reqBackend1 = new Counter('backend_1_requests');
export const reqBackend2 = new Counter('backend_2_requests');
export const reqBackend3 = new Counter('backend_3_requests');
export const errorRate = new Rate('custom_error_rate');
export const cpuLatency = new Trend('cpu_endpoint_duration', true);

export const options = {
    stages: [
        { duration: '1m', target: 100 },  // Tahap pemanasan (Warm-up) 100 VUs
        { duration: '3m', target: 500 },  // Beban menengah 500 VUs
        { duration: '4m', target: 1000 }, // Beban puncak (Peak saturation) 1.000 VUs
        { duration: '2m', target: 0 },    // Penurunan beban (Cool-down)
    ],
    thresholds: {
        http_req_failed: ['rate<0.05'], // Toleransi error rate maksimal 5%
        http_req_duration: ['p(95)<3000', 'p(99)<5000'], // SLO p95 < 3s, p99 < 5s
    },
};

const BASE_URL = __ENV.TARGET_URL || 'http://localhost';

export default function () {
    const res = http.get(`${BASE_URL}/api/cpu`, {
        timeout: '10s',
        tags: { endpoint: 'cpu_bound' },
    });

    const isSuccess = check(res, {
        'status is 200': (r) => r.status === 200,
        'has valid result': (r) => r.body && r.body.includes('success'),
    });

    errorRate.add(!isSuccess);
    cpuLatency.add(res.timings.duration);

    const backendServer = res.headers['X-Backend-Server'];
    if (backendServer) {
        if (backendServer.includes('web-backend-1')) {
            reqBackend1.add(1);
        } else if (backendServer.includes('web-backend-2')) {
            reqBackend2.add(1);
        } else if (backendServer.includes('web-backend-3')) {
            reqBackend3.add(1);
        }
    }

    sleep(0.1);
}
\end{lstlisting}

\newpage

% ==============================================================================
% LAMPIRAN 6: OTOMASI EKSEKUSI MATRIKS & AUDIT MAPE
% ==============================================================================
\section*{Lampiran 6: Skrip Otomasi Pengujian Matriks \& Analisis Presisi MAPE}

Berikut adalah skrip \textit{orchestration runner} (\texttt{run-matrix.sh}) untuk mengeksekusi 16 kombinasi eksperimen dan skrip Python (\texttt{parse-logs.py}) untuk menghitung metrik \textit{Mean Absolute Percentage Error} (MAPE):

\begin{lstlisting}[language=bash, caption={Skrip Otomasi 16 Skenario Matriks (run-matrix.sh)}]
#!/usr/bin/env bash
set -e

DIR="$(cd "$(dirname "${BASH_SOURCE[0]}")/.." >/dev/null 2>&1 && pwd)"
cd "$DIR"

mkdir -p logs/nginx logs/haproxy results

run_scenario() {
    local lb="$1"       # nginx / haproxy
    local algo="$2"     # rr / lc / wrr / wlc
    local workload="$3" # cpu / io

    echo ">>> MENJALANKAN SKENARIO: LB=$lb | ALGO=$algo | WORKLOAD=$workload <<<"
    rm -f logs/nginx/access.log logs/haproxy/haproxy.log

    if [ "$lb" == "nginx" ]; then
        export NGINX_CONF="./loadbalancer/nginx/nginx-${algo}.conf"
        docker compose --profile nginx down --volumes --remove-orphans > /dev/null 2>&1 || true
        docker compose --profile nginx up -d --build
        sleep 5
    else
        export HAPROXY_CONF="./loadbalancer/haproxy/haproxy-${algo}.cfg"
        docker compose --profile haproxy down --volumes --remove-orphans > /dev/null 2>&1 || true
        docker compose --profile haproxy up -d --build
        sleep 5
    fi

    local script="benchmark/${workload}-test.js"
    local output_json="results/${lb}_${algo}_${workload}.json"
    k6 run --out json="$output_json" "$script"

    local algo_type="weighted"
    if [ "$algo" == "rr" ] || [ "$algo" == "lc" ]; then
        algo_type="unweighted"
    fi

    if [ "$lb" == "nginx" ]; then
        python3 scripts/parse-logs.py --log logs/nginx/access.log --type nginx --algorithm "$algo_type"
    else
        python3 scripts/parse-logs.py --log logs/haproxy/haproxy.log --type haproxy --algorithm "$algo_type"
    fi
}

run_scenario "$1" "$2" "$3"
\end{lstlisting}

\vspace{0.4cm}

\begin{lstlisting}[language=Python, caption={Skrip Penghitung Deviasi Distribusi Beban MAPE (parse-logs.py)}]
#!/usr/bin/env python3
import sys, re, argparse
from collections import Counter

def parse_nginx_log(log_path):
    counts = Counter()
    pattern = re.compile(r'upstream_addr="([^"]+)"')
    total = 0
    with open(log_path, 'r', encoding='utf-8', errors='ignore') as f:
        for line in f:
            match = pattern.search(line)
            if match:
                addr = match.group(1).split(',')[0].strip()
                counts[addr] += 1
                total += 1
    return counts, total

def calculate_mape(actual_counts, is_weighted=True):
    if is_weighted:
        p_ideal = [100.0 * 1.0 / 7.0, 100.0 * 2.0 / 7.0, 100.0 * 4.0 / 7.0]
    else:
        p_ideal = [100.0 / 3.0, 100.0 / 3.0, 100.0 / 3.0]

    total_reqs = sum(actual_counts.values())
    if total_reqs == 0: return

    nodes_sorted = sorted(actual_counts.keys())
    deviations = []
    for idx, node in enumerate(nodes_sorted):
        vol = actual_counts[node]
        p_act = (vol / total_reqs) * 100.0
        p_id = p_ideal[idx] if idx < len(p_ideal) else (100.0 / len(nodes_sorted))
        dev = abs(p_act - p_id) / p_id * 100.0
        deviations.append(dev)

    mape = sum(deviations) / len(deviations) if deviations else 0.0
    print(f"Mean Absolute Percentage Error (MAPE): {mape:.2f}%")
\end{lstlisting}

\newpage

% ==============================================================================
% LAMPIRAN 7: SKRIP OTOMASI PENGUJIAN MATRIKS DAN TUNING KERNEL LINUX
% ==============================================================================
\section*{Lampiran 7: Skrip Otomasi Pengujian Matriks dan Tuning Kernel Linux}

Berikut adalah skrip otomasi eksekusi matriks pengujian beban dengan isolasi CPU core hibrida (\texttt{run-matrix.sh}) serta konfigurasi optimalisasi soket TCP kernel Linux (\texttt{99-skripsi.conf}):

\begin{lstlisting}[language=bash, caption={Skrip Otomasi Matriks Pengujian Beban dan Isolasi CPU (run-matrix.sh)}]
#!/usr/bin/env bash
# ==============================================================================
# AUTOMATED MATRIX TEST RUNNER (SKRIPSI AHMAD MUQTAFI)
# Lingkungan: Bare-Metal Laptop MSI Thin GF63 12UCX (Ubuntu 24.04 LTS)
# Isolasi CPU: Server pada P-Cores (0-7), k6 pada E-Cores (8-11)
# ==============================================================================
set -e

DIR="$(cd "$(dirname "${BASH_SOURCE[0]}")/.." >/dev/null 2>&1 && pwd)"
cd "$DIR"

mkdir -p logs/nginx logs/haproxy results

# Deteksi ketersediaan core CPU untuk taskset E-Cores (CPU 8-11)
TOTAL_CPUS=$(nproc 2>/dev/null || echo 1)
TASKSET_CMD=""
if [ "$TOTAL_CPUS" -ge 12 ]; then
    TASKSET_CMD="taskset -c 8-11"
    echo "[ISOLASI CPU] Grafana k6 dipin ke E-Cores (CPU 8-11) via taskset."
fi

run_scenario() {
    local lb="$1" algo="$2" workload="$3" quick_flag="$4"
    rm -f logs/nginx/access.log logs/haproxy/haproxy.log

    if [ "$lb" == "nginx" ]; then
        export NGINX_CONF="./loadbalancer/nginx/nginx-${algo}.conf"
        docker compose --profile nginx down --volumes --remove-orphans > /dev/null 2>&1 || true
        docker compose --profile nginx up -d --build
    else
        export HAPROXY_CONF="./loadbalancer/haproxy/haproxy-${algo}.cfg"
        docker compose --profile haproxy down --volumes --remove-orphans > /dev/null 2>&1 || true
        docker compose --profile haproxy up -d --build
    fi
    sleep 5

    local script="benchmark/${workload}-test.js"
    local output_json="results/${lb}_${algo}_${workload}_$(date +%Y%m%d_%H%M%S).json"
    local k6_extra_args=""
    [ "$quick_flag" == "--quick" ] && k6_extra_args="--vus 50 --duration 1m"

    if command -v k6 &> /dev/null; then
        $TASKSET_CMD k6 run $k6_extra_args --out json="$output_json" "$script"
    else
        docker run --rm -i --network="host" --cpuset-cpus="8-11" \
            -v "$DIR/benchmark:/benchmark" grafana/k6 run $k6_extra_args "/$script"
    fi
}
\end{lstlisting}

\vspace{0.4cm}

\begin{lstlisting}[language=bash, caption={Konfigurasi Tuning Parameter Kernel Linux (/etc/sysctl.d/99-skripsi.conf)}]
# Optimalisasi Batas Antrean Koneksi TCP dan Socket Buffer
net.core.somaxconn = 65535
net.ipv4.tcp_max_syn_backlog = 65535
net.ipv4.ip_local_port_range = 1024 65535
net.ipv4.tcp_tw_reuse = 1
net.ipv4.tcp_fin_timeout = 15
fs.file-max = 2097152
\end{lstlisting}
